\subsection{滤过群}

回忆（集合论，III，第145页），有序集G称为右（相应地，左）滤过的 \(^{1}\) ，如果对G的每一对元素 \((x, y)\) ，存在 \(z \in G\) 使得 \(x \leqslant z\) 且 \(y \leqslant z\) （相应地， \(z \leqslant x\) 且 \(z \leqslant y\) ）。每个右滤过的序群G也是左滤过的，反之亦然：事实上，由于存在 \(z \in G\) 使得 \(-x \leqslant z\) 且 \(-y \leqslant z\) ，我们就有 \(-z \leqslant x\) 且 \(-z \leqslant y\) （VI，第3页，推论）。因此我们将径直称之为滤过群。

命题4. --- 序群G是滤过的，当且仅当它由其正元素生成，也就是说，每个元素都是两个正元素之差。

确实，如果G是滤过的，那么对每个 \(x \in G\) 存在一个正元素 z，使得 \(x \leqslant z\) ，且 x 是两个正元素 z 与 z - x 的差。反之，若 x = u - v 且 y = w - t，其中 u、v、w、t 均为正元素，则元素 \(u + w\) 大于 x 且大于 y。

命题 5. --- 若 \((x_i)\) 是滤过群 G 的元素的有限族，则存在 \(z \in G\) 使得 \(x_i + z\) 对每个i都是正的。

如果 \(x_{i} = u_{i} - v_{i}\) ，其中 \(u_{i}\) 与 \(v_{i}\) 为正，只需取 \(z\) 为族 \((v_{i})\) 之和。

